\chapter{Fichier de configuration complet}
\label{ann:conf}

Le fichier ci-après est la configuration de référence livrée
avec le projet. Il fixe les valeurs identiques à celles utilisées
en interne avant l'introduction du système de configuration
externe, garantissant ainsi une parfaite compatibilité
descendante.

\inputminted[fontsize=\scriptsize, frame=lines, linenos, breaklines]{toml}{../aquarium.conf}

\chapter{Instructions de build et d'exécution}
\label{ann:build}

\section{Dépendances système}

\begin{itemize}
  \item GCC ou Clang (C11)
  \item \code{make} (GNU make)
  \item \code{pkg-config}
  \item GTK4 et ses headers de développement
  \item FFmpeg (pour le pipeline d'assets uniquement)
  \item Bear (optionnel, pour générer \file{compile\_commands.json})
\end{itemize}

\section{Commandes principales}

\begin{minted}{bash}
# Préparer les sprites à partir des fichiers bruts (à faire une fois)
make assets

# Compiler le projet
make

# Lancer l'aquarium
./build/bin/aquarium

# Nettoyer les artefacts de build
make clean

# Installer aquarium.conf dans ~/.config/gtkaqua/
make install-conf
\end{minted}

\section{Variables d'environnement utiles}

\code{XDG\_CONFIG\_HOME} contrôle où l'application cherche son
fichier de configuration utilisateur. Si cette variable n'est
pas définie, la valeur par défaut
\file{\textasciitilde/.config/} est utilisée.

\chapter{Arborescence du projet}
\label{ann:layout}

\begin{verbatim}
gtkaqua/
|-- Makefile
|-- README.md
|-- aquarium.conf            <- configuration externe
|-- include/
|   |-- config.h             <- variables globales extern
|   |-- conf.h               <- API du loader TOML
|   |-- entity.h
|   |-- species.h
|   |-- vec2.h               <- vecteurs 2D (header-only)
|   `-- world.h
|-- src/
|   |-- conf.c               <- loader TOML + validation
|   |-- entity.c             <- forces et integration physique
|   |-- main.c               <- GTK, UI, tick principal
|   |-- species.c            <- registre dynamique
|   `-- world.c              <- pool d'entites, orchestration
|-- vendor/
|   `-- tomlc99/             <- parseur TOML vendored
|       |-- toml.c
|       |-- toml.h
|       `-- LICENSE
|-- Fish/                    <- sprites bruts
|-- assets/                  <- sprites normalises + video de fond
|-- spritesheets/            <- references non utilisees
|-- docs/                    <- documentation projet
`-- build/
    |-- obj/                 <- fichiers .o
    `-- bin/aquarium         <- binaire final
\end{verbatim}

\chapter{Tableau récapitulatif des espèces}
\label{ann:species-table}

Le tableau~\ref{tab:species-summary} résume les principales
caractéristiques de chaque espèce telles que définies dans
\file{aquarium.conf}.

\begin{table}[h!]
\centering
\scriptsize
\setlength{\tabcolsep}{4pt}
\begin{tabular}{@{}lrrrrrrl@{}}
\toprule
\textbf{Espèce} & \textbf{Vit. max} & \textbf{Vit. min} &
\textbf{Force} & \textbf{Flock} & \textbf{Pop. max} & \textbf{Durée} &
\textbf{Proies} \\
\midrule
Goldfish   & 85  & 30 & 42 & 45 & 60 & 120 & -- \\
Clownfish  & 90  & 32 & 45 & 18 & 18 & 120 & -- \\
Angelfish  & 95  & 34 & 44 & 14 & 14 & 150 & -- \\
Barracuda  & 130 & 40 & 60 & 10 & 10 & 200 & goldfish, clownfish \\
Bass       & 112 & 36 & 53 & 10 & 10 & 180 & goldfish, clownfish \\
Trout      & 110 & 36 & 52 & 10 & 10 & 180 & goldfish \\
Tuna       & 125 & 40 & 56 &  4 &  8 & 200 & goldfish, clownfish \\
Shark      & 145 & 44 & 62 &  2 &  8 & 300 & tuna, bass, trout, angelfish, goldfish \\
\bottomrule
\end{tabular}
\caption{Récapitulatif des huit espèces par défaut.}
\label{tab:species-summary}
\end{table}

Les vitesses sont exprimées en pixels par seconde, la force
maximale en pixels par seconde au carré (accélération), et la
durée de vie en secondes.

\chapter{Définition de la structure \struct{SpeciesConfig}}
\label{ann:species-struct}

Pour référence, voici la définition complète du type qui porte
l'intégralité du paramétrage d'une espèce, telle qu'elle figure
dans \file{include/species.h}~:

\inputminted[fontsize=\footnotesize, frame=lines, linenos]{c}{../include/species.h}

\chapter{Signature de l'API de configuration}
\label{ann:conf-api}

L'interface publique exposée par le module de configuration~:

\inputminted[fontsize=\footnotesize, frame=lines, linenos]{c}{../include/conf.h}
